\section{Introduction}
\label{sec:introduction}
Real-time control software is increasingly deployed in domains where
timing faults may directly compromise safety, availability, or mission
success. In railway systems, timed automata have been used to model
train-gate controllers, communication protocols, and railway safety
layers \cite{basile2021unisig}. In automotive and autonomous-driving systems, they have
been used to reason about gear-control logic, road-junction behavior,
and time-bounded decision rules \cite{lindahl1998gear}. In medical cyber-physical
systems, timed models have been used to analyze pacemakers and
ventilator controllers, where physiological timing windows and actuation
deadlines are safety-critical \cite{jiang2012pacemaker}. In aerospace and avionics,
timed-automata models have supported schedulability and deadline
analysis for control software and distributed real-time architectures
\cite{david2010herschel}. Across these domains, correctness depends not only on the
order of discrete events, but also on when these events occur. A
controller that eventually raises an alarm, applies a pulse, sends a
message, or releases a resource may still be incorrect if the action
occurs too early, too late, or for the wrong duration.

Timed automata provide a well-established modeling formalism for this
class of systems \cite{alur1994theory}. Real-time systems are often represented as
networks of timed automata, where clocks, guards, invariants,
synchronizations, and resets describe both discrete behavior and
quantitative timing constraints \cite{bengtsson2004timed}. Verification is then performed
against temporal properties that may include safety, bounded response,
deadline, error-state reachability, alarm-duration, and deadlock-freedom
requirements. When a property is violated, a real-time model checker
such as UPPAAL, Kronos, or opaal can produce a timed diagnostic trace
(TDT), i.e., a timed counterexample that shows how the model reaches a
violating state \cite{daws1996kronos}.

However, diagnosis is only the first step. Interpreting a TDT and
turning it into a correct model edit remains a demanding engineering
task. Realistic control models often contain many locations,
transitions, clocks, guards, invariants, synchronizations, and resets. A
TDT exposes one feasible execution that violates a property, but it does
not directly identify which clock constraint is wrong, whether the
defect lies in a guard or an invariant, or how a bound should be changed
to restore the intended timing behavior \cite{koelbl2019clock}. This ambiguity becomes
more pronounced when the model is checked against multiple properties.
Control-system specifications are usually not single assertions, but
collections of timing requirements that jointly describe safety,
response, scheduling, and reliability obligations \cite{vogel2022patterns}. A local edit
that eliminates one diagnostic trace may still leave related traces
feasible, or may even make another property false. Thus, automated
repair for control models must go beyond explaining a single
counterexample and must validate each candidate repair against the
complete property set.

Existing timed-automata repair techniques provide an important
foundation. Clock-bound repair for timed systems encodes a TDT into
linear real arithmetic, introduces variation variables for clock bounds,
and uses MaxSMT solving to compute small modifications to guards and
invariants \cite{koelbl2019clock}. TarTar extends this idea into an automatic repair
analysis tool that suggests syntactic repairs and checks whether the
repaired model preserves the untimed functional behavior of the original
network timed automaton \cite{koelbl2020tartar}. Other related directions include
parameter synthesis for parametric timed automata \cite{alur1993parametric}, clock-guard
repair through abstraction and testing \cite{andre2019repairing}, robustness analysis for
uncertain timing constants \cite{bendik2021robustness}, and SAT/SMT-based model or program
repair \cite{jose2011bugassist}. These methods show that automated repair of timed
models is feasible, but most existing TDT-based repair workflows remain
centered on one violated property or one diagnostic trace at a time.

The engineering challenge addressed in this paper is different but
complementary to existing single-trace repair and synthesis techniques
\cite{andre2019repairing}. A straightforward repair loop can process one violated
property, collect one TDT, compute a local repair, instantiate a
candidate model, and then run regression verification. This loop is
useful as a baseline, but it may become inefficient or unstable when one
timing fault affects several properties or several related traces.
Similar traces may repeatedly encode overlapping parts of the same
faulty behavior; a candidate that eliminates one trace may still leave
other violations feasible; and a bound changed for one property may
affect another property when the same clock, transition, location, or
synchronization participates in both executions. If all editable bounds
are considered equally plausible repair targets, the solver may explore
many weakly related candidates, including edits that are syntactically
valid but distant from the timing decisions exposed by the observed
violations. Such repairs may pass a local check while giving engineers
little insight into the shared cause of the failures.

This paper presents \textbf{MPTA-Repair}, an automatic repair workflow
for multi-property, multi-counterexample clock-bound repair of
timed-automata control models. Given a faulty model, a set of timing
properties, and diagnostic traces produced by model checking,
MPTA-Repair suggests small modifications to numerical clock bounds in
transition guards and location invariants. The repair space is
intentionally restricted to clock-bound constants, rather than arbitrary
structural edits, specification changes, synchronization changes, or
reset changes. The workflow uses lightweight relation analysis among
properties, diagnostic traces, and editable clock bounds to guide repair
search, but this analysis is only an optimization strategy. Every
accepted candidate must satisfy the complete property set through
regression verification and must pass a TARTAR-style admissibility check
based on functional equivalence. In this way, MPTA-Repair aims to
generate repairs that are not only local counterexample fixes, but also
system-level timing corrections validated against the full set of
control requirements.

We evaluate MPTA-Repair on a suite of control-domain timed-automata
models covering railway alarm control \cite{gonzalez2024realtimedevs}, train-gate control
\cite{behrmann2004tutorial}, automotive gear control, dual-chamber pacemaker control,
autonomous road-junction rules, railway communication timeout handling,
mechanical ventilator control, and aerospace task scheduling. The
evaluation measures repair effectiveness, repair size, runtime,
regression results, and behavior preservation. We compare MPTA-Repair
with an iterative TARTAR-style baseline under the same full-property
regression and admissibility requirements, and report practical
observations from applying automated repair to industrial-style control
models.

This paper makes the following contributions:

\begin{enumerate}
\item
  We formulate clock-bound repair for timed-automata control models as a
  multi-property and multi-counterexample repair problem.
\item
  We present MPTA-Repair, an automatic repair workflow that uses
  diagnostic traces and lightweight relation analysis to guide
  clock-bound repair.
\item
  We implement a prototype repair pipeline combining diagnostic trace
  analysis, MaxSMT-based clock-bound repair, full-property regression
  verification, and TARTAR-style admissibility checking.
\item
  We evaluate the approach on a diverse control-domain benchmark suite
  and compare it with an iterative single-property repair baseline under
  the same acceptance criteria.
\item
  We report practical lessons on applying automated timed-automata
  repair to control models, especially regarding property selection,
  repair interpretability, and the need for full regression after each
  candidate repair.
\end{enumerate}

The remainder of this paper is organized as follows. Section II provides
the background and problem setting, including timed automata, diagnostic
traces, clock-bound repair, and the multi-property repair objective.
Section III presents MPTA-Repair and its relation-guided repair
workflow. Section IV describes the prototype implementation and
evaluation setup. Section V reports the evaluation results and practical
lessons learned from the control-domain models. Section VI discusses
threats to validity and related work. Section VII concludes the paper.
